\chapter{Marco Teórico y Estado del Arte}
\label{ch:marco_teorico}

Este capítulo presenta los fundamentos sobre los que se construye la investigación: principios neurofisiológicos del EEG, bases matemáticas del Procesamiento de Señales en Grafos, estado del arte en interpolación de señales biomédicas, y los mecanismos internos que convierten a la implementación de MNE-Python en una referencia comunitaria sólida.

\section{Fundamentos Neurofisiológicos de la Señal EEG}

\subsection{Generación y Propagación}

La señal de voltaje EEG registrada en el cuero cabelludo resulta de la actividad sincronizada de grandes poblaciones de neuronas piramidales en la corteza cerebral \cite{buzsaki2006rhythms}. Estas neuronas, orientadas perpendicularmente a la superficie cortical, generan dipolos de corriente ante potenciales postsinápticos. La suma espacial de estos dipolos ---se requieren entre diez mil y cien mil neuronas sincronizadas para producir una señal detectable \cite{nunez2006electric}--- constituye la fuente primaria del EEG.

La señal de voltaje generada en la corteza atraviesa desde su origen hasta los electrodos de superficie múltiples capas de conductividad variable: líquido cefalorraquídeo, meninges, cráneo (aproximadamente ochenta veces menos conductor que el cuero cabelludo) y cuero cabelludo \cite{nunez2006electric}. Esta conducción volumétrica actúa como un filtro espacial paso-bajo que suaviza la distribución de potencial y limita la resolución espacial efectiva del EEG a unos dos o tres centímetros. La Figura~\ref{fig:origen_eeg}(a) ilustra este proceso de propagación desde la corteza hasta el cuero cabelludo, y la Figura~\ref{fig:origen_eeg}(b) muestra la configuración del campo eléctrico generado por un dipolo cortical.
\begin{figure}[htbp]
\centering
\includegraphics[width=0.94\textwidth]{capitulo2_origen_eeg}
\caption[Generación de la señal EEG y manifestaciones relevantes]{Generación de la señal EEG y manifestaciones relevantes para la interpolación de canales. (a) Propagación desde fuentes corticales hasta electrodos de superficie a través de capas con conductividad distinta, representadas sobre una referencia de cabeza. (b) Campo idealizado de un dipolo cortical y sus líneas de potencial. (c) Ejemplo esquemático de ERP con componentes N100, P200 y P300, cuya morfología debe preservarse al interpolar. (d) Topografía de electrodos con canales observados y canales fallados sobre una cabeza con orientación anatómica. Esquema de elaboración propia basado en la descripción neurofísica de \cite{nunez2006electric} y en la nomenclatura ERP de \cite{luck2014introduction}.}
\label{fig:origen_eeg}
\end{figure}
\subsection{Bandas de Frecuencia y Potenciales Relacionados a Eventos}

La señal EEG se caracteriza por su contenido oscilatorio en bandas de frecuencia asociadas a estados funcionales específicos: delta (0.5--4~Hz, sueño profundo), theta (4--8~Hz, procesamiento de memoria), alpha (8--13~Hz, reposo vigil), beta (13--30~Hz, procesamiento activo) y gamma (30--100~Hz, procesamiento perceptual de alto nivel) \cite{buzsaki2006rhythms}. La preservación de la estructura espectral en cada banda es un criterio fundamental para evaluar la calidad de cualquier interpolación.

Los Potenciales Relacionados a Eventos (ERP) son fluctuaciones de voltaje sincronizadas con un estímulo sensorial, motor o cognitivo, obtenidas mediante promediación coherente de múltiples ensayos \cite{luck2014introduction}. Componentes como N100, P300 y N400 constituyen biomarcadores electrofisiológicos cuya preservación morfológica ---amplitud de pico, latencia y distribución topográfica--- es uno de los criterios de evaluación más exigentes para cualquier método de interpolación. La Figura~\ref{fig:origen_eeg}(c) muestra un ejemplo de forma de onda ERP con sus componentes principales, y la Figura~\ref{fig:origen_eeg}(d) ilustra la topografía de canales buenos y fallados.

\section{Procesamiento de Señales en Grafos}

El Procesamiento de Señales en Grafos (GSP) \cite{shuman2013,ortega2018} extiende los conceptos del procesamiento clásico de señales a dominios irregulares representados por grafos.

\subsection{Definiciones Fundamentales}

Un grafo $\mathcal{G} = (\mathcal{V}, \mathcal{E}, \mathbf{W})$ consiste en $N$ nodos $\mathcal{V} = \{v_1, \dots, v_N\}$ ---los electrodos---, un conjunto de aristas $\mathcal{E} \subseteq \mathcal{V} \times \mathcal{V}$ que define qué electrodos están conectados, y una matriz de pesos $\mathbf{W} \in \mathbb{R}^{N \times N}$ simétrica que especifica con qué intensidad: $W_{ij} > 0$ si existe arista entre $i$ y $j$, y $W_{ij} = 0$ en caso contrario. La distinción entre $\mathcal{E}$ y $\mathbf{W}$ es relevante porque $\mathcal{E}$ define la topología (pares conectados) mientras que $\mathbf{W}$ codifica la fuerza de cada conexión; en grafos densos pueden coincidir, pero en grafos sparse como k-NN difieren.

Una señal de grafo en el instante $t$ es el vector $\mathbf{x}_t \in \mathbb{R}^N$ de voltajes en los $N$ electrodos. La matriz $\mathbf{X} \in \mathbb{R}^{N \times T}$ contiene $T$ instantes temporales de la ventana de análisis EEG.

\subsection{Laplaciano, Suavidad y Transformada de Fourier en Grafos}

La matriz Laplaciana se define como $\mathbf{L} = \mathbf{D} - \mathbf{W}$, con $D_{ii} = \sum_j W_{ij}$. La forma cuadrática
\begin{equation}
S(\mathbf{x}) = \mathbf{x}^\top \mathbf{L} \mathbf{x} = \frac{1}{2} \sum_{i,j} W_{ij} (x_i - x_j)^2
\label{eq:laplacian_quadratic}
\end{equation}
cuantifica la variación total del vector de voltajes $\mathbf{x} \in \mathbb{R}^N$ sobre el grafo: valores pequeños indican que nodos fuertemente conectados tienen valores similares ---la señal es \textit{suave}---. Esta propiedad es el pilar de la interpolación GSP: la mejor reconstrucción minimiza $S(\mathbf{x})$ sujeta a coincidir con los canales observados. La Ecuación~\ref{eq:laplacian_quadratic} es la base de todos los métodos de regularización en grafos discutidos en este capítulo.

$\mathbf{L}$ admite la descomposición espectral $\mathbf{L} = \mathbf{U} \boldsymbol{\Lambda} \mathbf{U}^\top$, donde $\boldsymbol{\Lambda} = \operatorname{diag}(\lambda_1, \dots, \lambda_N)$ contiene los autovalores ordenados $0 = \lambda_1 \leq \cdots \leq \lambda_N$, y $\mathbf{U}$ es la matriz ortogonal de autovectores. Los $\lambda_k$ actúan como frecuencias en el grafo: valores pequeños corresponden a modos de variación lenta, valores grandes a modos de variación rápida (Figura~\ref{fig:gsp_conceptos}b).

La Transformada de Fourier en Grafos (GFT) se define como $\hat{\mathbf{x}} = \mathbf{U}^\top \mathbf{x}$. Una señal es \textit{de banda limitada} si $\hat{x}_k = 0$ para $k > K$, con $K \ll N$. La suavidad en el dominio espectral es $S(\mathbf{x}) = \sum_k \lambda_k \hat{x}_k^2$, penalizando los modos de alta frecuencia proporcionalmente a su autovalor. En este contexto, una reconstrucción regularizada es la estimación de los canales faltantes que conserva los canales observados y minimiza simultáneamente una penalización de variación sobre el grafo.

\begin{figure}[htbp]
\centering
\includegraphics[width=0.96\textwidth]{ch2_gsp_concepts_clean}
\caption[Conceptos fundamentales de GSP]{Conceptos fundamentales de GSP aplicados a EEG. (a) Los electrodos forman un grafo ponderado sobre una topografía de cabeza; las aristas codifican cercanía espacial. (b) El espectro del Laplaciano ordena los modos desde variación suave hasta variación rápida. La curva azul muestra autovalores, mientras que la curva naranja ilustra que la energía de una señal suave tiende a concentrarse en modos bajos. (c) La reconstrucción regularizada estima los canales faltantes combinando fidelidad a los canales observados y suavidad inducida por el grafo.}
\label{fig:gsp_conceptos}
\end{figure}

La Figura~\ref{fig:gsp_conceptos}(a) muestra el grafo de electrodos, la Figura~\ref{fig:gsp_conceptos}(b) separa explícitamente los autovalores del Laplaciano de la energía modal ilustrativa, y la Figura~\ref{fig:gsp_conceptos}(c) resume la reconstrucción regularizada como compromiso entre datos observados y suavidad espacial.

\subsection{Construcción de Grafos para EEG}
\label{sec:graph_construction}

La elección del grafo subyacente es una decisión metodológica de primer orden: un mismo algoritmo de interpolación GSP produce resultados distintos según el grafo sobre el que opera. En este trabajo se implementaron y evaluaron diez estrategias de construcción de grafos, organizadas en tres categorías según la fuente de información que utilizan.

\textbf{Grafos geométricos (basados en distancia).} Utilizan exclusivamente las posiciones físicas de los electrodos, lo que los hace inmunes a la presencia de canales faltantes y deterministas por construcción:

\begin{itemize}[leftmargin=*,itemsep=1pt,topsep=2pt]
\item \textbf{k-NN no ponderado (knn).} Conecta cada electrodo a sus $k$ vecinos más próximos con pesos binarios $\{0,1\}$. Es el grafo más simple y sirve como línea base estructural.
\item \textbf{k-NN gaussiano (knng).} Extiende el anterior ponderando cada arista con $W_{ij} = \exp(-\|\mathbf{p}_i - \mathbf{p}_j\|_2^2 / 2\sigma^2)$, donde $\mathbf{p}_i \in \mathbb{R}^3$ es el vector de coordenadas 3D del electrodo $i$ según el sistema 10--20 o 10--10, $\sigma$ controla la escala de decaimiento y $k$ determina la dispersión.
\item \textbf{k-NN gaussiano variable (vknng).} Similar al knng pero con $k_i$ adaptativo según la densidad local de electrodos: en regiones densas (e.g., zona occipital en montajes 10--10) usa menos vecinos, y en regiones dispersas usa más. El parámetro $\alpha$ controla la sensibilidad a la densidad local.
\item \textbf{$\varepsilon$-ball (epsilon\_ball).} Conecta todos los pares cuya distancia es menor que un umbral $\varepsilon$, con pesos $\{0,1\}$. A diferencia de k-NN, el grado de cada nodo es variable y depende de la densidad local.
\item \textbf{Distancia completa gaussiana (gaussian).} Conecta todos los pares de electrodos con el núcleo gaussiano, generando una matriz densa. No requiere el parámetro $k$ y captura interacciones de largo alcance, aunque puede introducir ruido de conexiones débiles.
\item \textbf{Conectividad total inversa (fully\_connected\_inverse\_distance).} Usa $W_{ij} = 1/d_{ij}$, generando un grafo denso donde la influencia decae linealmente con la distancia.
\end{itemize}

\textbf{Grafos basados en correlación funcional.} Utilizan la serie temporal observada para cuantificar similitud entre electrodos. La construcción es $W_{ij} = |\rho_{ij}|$, donde $\rho_{ij}$ es el coeficiente de correlación de Pearson entre los electrodos $i$ y $j$, calculado solo con canales disponibles. Este criterio es informativo para exploración, pero introduce circularidad si se usa para confirmar desempeño: el grafo queda parcialmente determinado por las mismas señales sobre las que luego se evalúa la reconstrucción. Por esta razón, los grafos de correlación se usaron solo en la fase de selección preliminar y no sostienen las conclusiones comparativas principales.
\textbf{Grafos data-driven (aprendidos desde los datos).} No requieren especificar un núcleo ni sus parámetros; la estructura del grafo emerge de las señales mismas:

\begin{itemize}[leftmargin=*,itemsep=1pt,topsep=2pt]
\item \textbf{Kalofolias (kalofolias).} Formula el aprendizaje del grafo como optimización convexa que minimiza la variación de la señal observada con restricciones de grado positivo \cite{kalofolias_how_2016}. Se implementó una versión propia con descenso de gradiente proyectado sobre la función objetivo de log-grados más regularización de Frobenius.
\item \textbf{NNK (nnk).} Regresión de kernel no-negativa que selecciona automáticamente los vecinos relevantes para cada nodo resolviendo un problema NNLS (\textit{Non-Negative Least Squares}) local, sin necesidad de especificar $k$ \cite{shekkizhar2020}.
\item \textbf{AEW (Adaptive Edge Weighting).} Optimización conjunta de pesos y topología mediante el algoritmo AEW \cite{shekkizhar2020}, que alterna entre actualización de pesos vía descenso de gradiente y proyección sobre restricciones de grado.
\end{itemize}
Los grafos data-driven se evaluaron en la fase exploratoria pero no se incluyeron en la fase comparativa por dos razones: su mayor costo computacional ($O(N^3)$ en el caso de Kalofolias) y la dificultad de garantizar reproducibilidad determinista entre ejecuciones. Esta decisión no es puramente teórica: el Capítulo~\ref{ch:resultados} documenta la selección preliminar de combinaciones método--grafo y la reducción experimental que llevó a concentrar la comparación final en grafos geométricos reproducibles y en la referencia MNE-Python. Su impacto sobre TRSS se propone como trabajo futuro.

\section{Estado del Arte en Interpolación de Señales Biomédicas}

\subsection{Métodos Geométricos y Estadísticos}

El método más establecido es el de \textit{spherical splines} \cite{perrin1989}, que modela el potencial sobre el cuero cabelludo como una función suave sobre una esfera. El potencial estimado en una posición $\mathbf{r}$ es $\hat{V}(\mathbf{r}) = c_0 + \sum_{i=1}^{M} c_i \, g_m(\mathbf{r}, \mathbf{r}_i)$, donde $M$ es el número de canales observados, $g_m$ es la función de Green del operador pseudo-laplaciano de orden $m$ sobre la esfera y los coeficientes se obtienen resolviendo un sistema lineal aumentado. En una forma regularizada, el bloque principal del sistema puede escribirse como $(\mathbf{G} + \lambda \mathbf{I})\mathbf{c} + c_0\mathbf{1}=\mathbf{v}$, donde $G_{ij}=g_m(\mathbf{r}_i,\mathbf{r}_j)$ y $\lambda$ estabiliza el ajuste. Así, $m$ controla la rigidez de la superficie y $\lambda$ regula el compromiso entre ajuste y suavidad.
Junto a los splines esféricos, se implementaron y evaluaron otros métodos de la familia geométrica. La \textbf{interpolación por distancia inversa} (IDW) estima el canal faltante como promedio ponderado de los canales observados, con pesos $w_i = 1/d_i^p$ donde $d_i$ es la distancia euclidiana al canal faltante y $p$ controla la velocidad de decaimiento. La \textbf{interpolación por vecino más cercano} asigna al canal faltante el valor del canal observado más próximo. Las \textbf{funciones de base radial} (RBF) generalizan el enfoque usando núcleos gaussianos o multicuádricos y han sido aplicadas explícitamente a reconstrucción de EEG \cite{jager2016rbf}. La \textbf{interpolación lineal} estima por triangulación sobre las posiciones 2D/3D de los electrodos. Estos métodos comparten la propiedad de operar exclusivamente en el dominio espacial, ignorando la estructura temporal de la señal.

Entre los métodos estadísticos, el \textbf{Análisis de Componentes Independientes (ICA)} \cite{jung2001imaging} se utilizó en dos variantes: una implementación propia con FastICA y la implementación de MNE-Python (ICA-MNE) con algoritmo Picard. ICA proyecta las señales observadas sobre fuentes independientes, elimina componentes artefactuales y reconstruye los canales faltantes desde el espacio de fuentes. Su inclusión se justifica porque es el método estándar de la comunidad EEG para remoción de artefactos, y evaluar su capacidad de reconstruir canales completamente ausentes —no solo corregir canales ruidosos— es una pregunta metodológica relevante. Sin embargo, ICA requiere la estimación de la matriz de mezcla sobre datos completos y su desempeño se degrada cuando la pérdida de canales es severa y contigua, por lo que se reservó para la fase exploratoria. El \textbf{Análisis de Componentes Principales (PCA)} se menciona aquí como contexto histórico pero no se evaluó experimentalmente por su desempeño consistentemente inferior en pruebas preliminares.
\subsection{Métodos Basados en GSP}

La \textbf{regularización de Tikhonov en grafos} formula la interpolación como un problema de optimización convexa:
\begin{equation}
\hat{\mathbf{x}} = \arg\min_{\mathbf{x}} \; \|\mathbf{M} \odot (\mathbf{x} - \mathbf{y})\|_2^2 \;+\; \alpha \cdot \mathbf{x}^\top \mathbf{L} \mathbf{x},
\label{eq:tikhonov_graph}
\end{equation}
donde $\mathbf{M} \in \{0,1\}^N$ es la máscara binaria de canales observados (con $M_i = 1$ si el canal $i$ está observado y $M_i = 0$ si está faltante), $\odot$ denota el producto Hadamard (multiplicación elemento a elemento entre dos matrices o vectores del mismo tamaño), y $\alpha$ controla el compromiso entre fidelidad a los datos y suavidad espacial. La Ecuación~\ref{eq:tikhonov_graph} admite solución cerrada $\hat{\mathbf{x}} = (\operatorname{diag}(\mathbf{M}) + \alpha \mathbf{L})^{-1} (\mathbf{M} \odot \mathbf{y})$. El \textbf{método de muestreo en grafos} asume que la señal es de banda limitada y busca la señal de banda limitada que interpola exactamente los valores observados.

\subsection{Regularización Espacio-Temporal: TRSS}

La Regularización Temporal de Suavidad de Sobolev (TRSS) \cite{giraldo2022} incorpora la continuidad temporal de las señales neuronales mediante un término adicional de penalización. El problema conjunto sobre la matriz espacio-temporal $\mathbf{X} \in \mathbb{R}^{N \times T}$ es
\begin{equation}
\hat{\mathbf{X}} = \arg\min_{\mathbf{X}} \; \|\mathbf{M} \odot (\mathbf{X} - \mathbf{Y})\|_F^2 \;+\; \alpha \cdot \operatorname{tr}(\mathbf{X}^\top \mathbf{L} \mathbf{X}) \;+\; \beta \cdot \|\mathbf{X} \mathbf{D}_t\|_F^2,
\label{eq:trss}
\end{equation}
donde $\mathbf{D}_t \in \mathbb{R}^{T \times (T-1)}$ es la matriz de diferencias finitas con $(\mathbf{D}_t)_{t,t} = -1$, $(\mathbf{D}_t)_{t+1,t} = 1$, y $\|\mathbf{X} \mathbf{D}_t\|_F^2 = \sum_{i=1}^N \sum_{t=1}^{T-1} (X_{i,t+1} - X_{i,t})^2$ penaliza los cambios abruptos muestra a muestra. $\beta$ controla la suavidad temporal: $\beta = 0$ reduce TRSS a Tikhonov independiente por instante; $\alpha = 0$ lo convierte en un suavizador temporal sin guía espacial. La interacción de ambos términos se evalúa empíricamente mediante el estudio de descomposición de componentes del Capítulo~4. La Figura~\ref{fig:trss_operator} descompone visualmente la contribución de cada término en la Ecuación~\ref{eq:trss}.

\begin{figure}[htbp]
\centering
\includegraphics[width=0.94\textwidth]{ch2_trss_operator}
\caption[Descomposición del operador TRSS]{Descomposición del operador TRSS. La reconstrucción combina fidelidad a los canales observados, suavidad espacial inducida por el Laplaciano del grafo y suavidad temporal inducida por diferencias finitas. El término temporal no agrega información externa: restringe el espacio de soluciones admisibles mediante $D_t$.}
\label{fig:trss_operator}
\end{figure}

\section{Mecanismo Interno de la Interpolación en MNE-Python}
\label{sec:mne_interpolate_bads}

Esta sección documenta las mejoras de ingeniería ---no descritas en la literatura académica--- que convierten a la implementación de \textit{spherical splines} en MNE-Python \cite{gramfort2013meg,gramfort2014mne} en una referencia sólida. La información se obtuvo mediante inspección directa del código fuente.

\subsection{Formulación Original y sus Seis Mejoras}

El artículo de Perrin \cite{perrin1989} formula la interpolación como un sistema lineal aumentado en el que la regularización se incorpora al bloque de Green, de manera análoga a $(\mathbf{G}+\lambda\mathbf{I})\mathbf{c}$. La formulación original ya utiliza todos los canales observados para estimar cada canal faltante; por tanto, la ventaja de MNE-Python no proviene de usar más canales que Perrin, sino de una implementación numéricamente robusta y bien integrada al flujo M/EEG. MNE-Python incorpora las siguientes decisiones de ingeniería:
\begin{enumerate}[leftmargin=*,itemsep=2pt,topsep=4pt]
\item \textbf{Normalización a esfera unitaria.} Las coordenadas de electrodos $\mathbf{p}_i$ se centran y escalan: $\mathbf{p}_i' = (\mathbf{p}_i - \Bar{\mathbf{p}}) / \max_j \|\mathbf{p}_j - \Bar{\mathbf{p}}\|_2$. Esto elimina la dependencia de unidades físicas y mejora el condicionamiento numérico.
\item \textbf{Uso global preservado.} Igual que la formulación de Perrin, todos los $M$ canales observados contribuyen a cada estimación. En un montaje de 64 canales con 2 canales malos, la reconstrucción aprovecha 62 electrodos; esto la diferencia de métodos locales como vecino más cercano o k-NN con $k$ pequeño, no de Perrin.
\item \textbf{Regularización adaptativa.} El regularizador efectivo se escala según la matriz de Green $\mathbf{G}$: $\varepsilon = 10^{-6} \cdot \operatorname{tr}(\mathbf{G}) / M$. Aquí $\varepsilon$ cumple el papel numérico de $\lambda$ en el sistema regularizado.
\item \textbf{Aumentación con restricciones lineales.} Se fuerzan $\sum_i c_i = 0$ y condiciones análogas para las coordenadas, garantizando unicidad.
\item \textbf{Pseudoinversa con tolerancia adaptativa.} Ante matrices numéricamente singulares, se usa una pseudoinversa con umbral proporcional a la norma.
\item \textbf{Reintento automático.} Si la reconstrucción produce valores nulos, el sistema reintenta con parámetros por defecto.
\end{enumerate}

Estas mejoras transforman un algoritmo teóricamente simple en un estimador notablemente robusto. La implicación para esta tesis es clara: no es válido comparar contra implementaciones simplificadas de Perrin; la referencia debe ser MNE-Python mismo.

\subsection{Los Tres Pilares de la Robustez}

Desde el punto de vista \textbf{físico}, el modelo de cabeza esférica ---cuya solución involucra polinomios de Legendre como base ortogonal natural para la ecuación de Laplace en coordenadas esféricas--- es una solución regularizada de dicha ecuación: no es una heurística sino el modelo más simple que captura la física de la conducción volumétrica. Cuando la información previa es correcta ---y la suavidad espacial del EEG es un hecho físico bien documentado---, los métodos con restricciones físicas simples superan a alternativas más complejas.

Desde el punto de vista \textbf{estadístico}, la regularización de MNE puede interpretarse como un estimador penalizado con una preferencia explícita por superficies suaves. No se afirma que MNE implemente un modelo bayesiano completo; la analogía relevante es que la penalización actúa como información previa de suavidad espacial. En el régimen de alta redundancia espacial típico del EEG de cuero cabelludo, este estimador penalizado aprovecha una propiedad física real de la señal.
Desde el punto de vista \textbf{informacional}, el EEG de cuero cabelludo exhibe redundancia espacial alta por el filtrado paso-bajo de la conducción volumétrica. Los métodos globales, que utilizan todos los canales observados, extraen más información que los métodos locales que restringen la estimación a un subconjunto de vecinos, como el vecino más cercano o k-NN con $k$ pequeño.
La Tabla~\ref{tab:method_families} resume las familias de métodos discutidas en esta sección, indicando su principio matemático, fortaleza principal y limitación característica.

\begin{table}[htbp]
\centering
\caption{Familias de métodos: principio matemático, fortalezas y limitaciones.}
\label{tab:method_families}
\small
\begin{tabular}{@{}p{2.2cm}p{3.0cm}p{3.1cm}p{3.2cm}@{}}
\toprule
Familia & Principio & Fortaleza & Limitación \\
\midrule
Geométrica & Distancia + superficie esférica & Velocidad, simplicidad & Ignora dinámica temporal \\
Estadística & Descomposición en fuentes (ICA) & Separa artefactos & Requiere datos completos \\
GSP-Tikhonov & Suavidad espacial en grafo & Topología flexible & Independiente por instante \\
GSP-Muestreo & Banda limitada en grafo & Fundamento de muestreo & Sensible a selección de $K$ \\
TRSS & Suavidad espacio-temporal & Preserva morfología & Costo computacional alto \\
MNE-Python & Splines + mejoras ing. & Robusta en rég. favorable & Pierde precisión bajo pérdida contigua \\
\bottomrule
\end{tabular}
\end{table}